% !TeX root = ../main.tex
\chapter{Literature Review}
\label{chap:lit_review}

\subsection{Definition}
\label{subsec:lit:definition}
Biomedical \gls{ie} is the process of automatically identifying structured information - entities such as genes, diseases, chemicals and species, and the relationships between them - from unstructured biomedical text such as journal abstracts and full-test articles. This process is usually comprised of two main tasks: \gls{ner}, which identifies biomedical entities, and \gls{re}, which identifies the relationships between those entities \cite{perera2020_ie_introduction}. ------

\subsection{Motivation}
\label{subsec:lit:motivation}
The volume of biomedical literature has expanded at a rate that exceeds the capacity of manual review, with previous noncurrent analyses demonstrating an exponential growth in the number of biomedical publications \cite{lu2011_pubmed_publicationrate}. This growth is supported by more recent data - a mapping of over 20 million PubMed abstracts has demonstrated the magnitude of contemporary biomedical literature to the level that the volume of biomedical and life-science publications has made it difficult for researchers to keep track of new literature \cite{gonzalez2024_biomed_landscape}. This is not only a volume issue, but also an issue with discovery - log analysis of PubMed search behaviour shows that over 80\% of views fall within the first page (top 20 results) of returned literature \cite{islamaj2009_firstpage}, meaning a substantial majority of literature goes largely unseen by researchers relying on manual review of ranked results.

\subsection{Tasks}
\label{subsec:lit:tasks}

\subsection{Challenges}
\label{subsec:lit:challenges}

\subsection{Traditional Methods}
\label{subsec:lit:traditional_methods}

\subsection{Transition to LLMs}
\label{subsec:lit:transition_to_llms}